% Vista pública derivada de 73_rev2_electron_lectores_torres.tex: cuerpo exacto y único retítulo académico del lector temporal de 108 estados. El fragmento del handoff permanece byte-exacto.
\section{Construcción completa del electrón beta: ruta central, firma, carácter y masa rectora}
\Needspace{7\baselineskip}
\subsection{Ruta central del electrón y secuencia de registros genealógicos}
El electrón ocupa el único centro \((5,5)\), con orientación activa \(++\):

\[
 \Gamma_e=(5,5,++),\qquad
 \tau_{12}=+++---+++---.
\]
Sus lectores de desplazamiento son

\[
 (D_1,D_3,D_4,D_{27},D_{36})=(54,56,68,48,32).
\]
No deben fundirse:

\begin{enumerate}
\item la firma cruda \((24,0,12,0,-12)\);
\item la firma par CT--108 \((-12,-4,12,-6,12)\);
\item el vector de doce eventos;
\item el origen afín \(v_e=0\) del carácter relativo.
\end{enumerate}
Tampoco debe confundirse \(\Gamma_e\) con la palabra \(w_e=201101\) del lector de la constante de Euler \(e\).

\Needspace{7\baselineskip}
\subsection{Selector basal del electrón y transición al valor beta rector}
El selector interno, construido sin usar la masa observada, cierra

\[
 -22=-\left|((\mathbb Z/9\mathbb Z)\times\mathbb F_3)\setminus\iota(R_\pi)\right|,
 \qquad
 +15=|R_\pi\times\mathbb F_3|.
\]
La lectura de vacancias aporta además

\[
 V_e=\begin{pmatrix}21&22\\6&7\end{pmatrix},\qquad
 \det V_e=15.
\]
El signo negativo de 22 registra déficit/corte; el positivo de 15, retención/área orientada. La matriz es local al electrón y no se extiende como selector universal.

La etapa basal es

\[
 e_0=\frac{\sqrt3}{4}
 \left[\exp\!\left(\frac{\varphi}{\pi^2}\right)
 -22\alpha_{\rm HMT}^3\right](1+15\Delta_4),
\]
\[
 e_0=0.510998922488\ldots\ \mathrm{MeV}.
\]
Las genealogías enteras son

\[
 22=\frac{396}{18}=27-5=24-2,
 \qquad
 15=\frac{270}{18}=\binom62=\binom64.
\]
\Needspace{7\baselineskip}
\subsection{Lectores bidireccionales y beta}
Para toda ruta,

\[
 K_D(\Gamma)=\left(D_{27}+D_{36},D_1,\frac{D_4-D_3}{2}\right),
\]
\[
 K_\Omega(\Gamma)=(\Omega_{90,120},54,P_6).
\]
Sobre las 324 rutas aparecen 14 perfiles \(K_D\), 9 perfiles \(K_\Omega\), 40 pares conjuntos y 18 coincidencias. El único valor común es

\[
 (80,54,6).
\]
Este 80 procede del lector bidireccional de la ruta electrónica, no del período 80 de la palabra de Euler. Por ello

\[
 \kappa_e=80-54\alpha_{\rm HMT}+6\Delta_4,
\]
\[
 \boxed{
 \mathfrak e_e^\beta=e_0R_{\rm act}^{\kappa_e}
 =0.5109989506900008086082571648\ldots\ \mathrm{MeV}.}
\]
Esta es la salida electrónica rectora. El basal no la sustituye. La igualdad con la masa de la antipartícula procede de la conjugación de ruta y de la paridad del operador de masa, no de añadir una segunda fila numérica.

\Needspace{7\baselineskip}
\section{Realización física por acoplamiento espectral}
\Needspace{7\baselineskip}
\subsection{Obstrucción a la selección puntual equivariante de una multisección no central}
Sea \(F\subset\mathcal R_{324}\) una fibra de rutas compatible con una clase estructural. Su espacio lineal es

\[
 \mathcal H_F=\bigoplus_{\Gamma\in F}\mathbb C e_\Gamma.
\]
El carácter central y los lectores definen operadores diagonales:

\[
 \widehat{\mathcal A}_F^{(0)}e_\Gamma
 =\chi_{\rm full}(\sigma_\Gamma,\eta_\Gamma)e_\Gamma,
\]
\[
 \widehat K_F^r e_\Gamma=\kappa_r(\Gamma)e_\Gamma,
 \qquad r\in\{D,\Omega\}.
\]
La corrección two--way es

\[
 \widehat{\mathcal M}_F^{\beta,r}
 =R_{\rm act}^{\widehat K_F^r/2}
  \widehat{\mathcal M}_F^{(0)}
  R_{\rm act}^{\widehat K_F^r/2}.
\]
Cuando \(\dim\mathcal H_F>1\), estos operadores tienen, en general, varios autovalores. Elegir uno porque se aproxima a una masa conocida invertiría la causalidad. Tampoco es correcto exigir que un nombre físico señale una celda única: la especie no central reside en una representación de la fibra.

\Needspace{7\baselineskip}
\subsection{Funtor espectral de la realización física: núcleo positivo de acoplamiento, subespacio persistente y medida espectral}
Una preparación física \(P\) aporta datos de representación, no una masa:

\[
 \mathfrak r(P)=
 (q,\,j,\,\epsilon_{2\pi},\,\epsilon_{4\pi},\,
  \text{estadística},\,\text{generación},\,\text{sabor},\,
  \text{órbita de color},\,\text{canal de interacción}).
\]
Sea \(\mathfrak M\) un acoplamiento compatible con esa representación. Se define un núcleo positivo

\[
 K_{P,\mathfrak M}:F\times F\longrightarrow\mathbb C
\]
que satisface:

\begin{enumerate}
\item \(K=K^*\) y \(K\ge0\);
\item \(K(\Gamma,\Gamma')=0\) si las rutas violan carga, paridad, color, estadística o
canal de \(P\);

\item \(K\) conmuta con las simetrías que no rompe el acoplamiento;
\item \(K(C_M\Gamma,C_M\Gamma')=\overline{K(\Gamma,\Gamma')}\) bajo conjugación;
\item sus entradas se calculan con registro genealógico, frontera, memoria, acarreo y holonomía, sin
magnitudes externas.

\end{enumerate}
El operador de evolución acoplado es

\[
 \mathcal U_{P,\mathfrak M}
 =K_{P,\mathfrak M}^{1/2}\,\mathcal U_{\rm TPK}\,
  K_{P,\mathfrak M}^{1/2}.
\]
El subespacio persistente se obtiene mediante la proyección espectral sobre los modos que sobreviven al retorno y a la poda:

\[
 \Pi_{P,\mathfrak M}^{\rm surv}
 =\mathbf1_{\mathcal S_{\rm surv}}(\mathcal U_{P,\mathfrak M}).
\]
La realización física correcta es el funtor espectral

\[
 \boxed{
 \mathcal S_{\rm phys}^{\rm spec}(P,\mathfrak M)
 =\left(
 \mathcal H_{P,\mathfrak M}^{\rm surv},
 \widehat{\mathcal M}_{P,\mathfrak M},
 \mu_{P,\mathfrak M}
 \right),}
\]
donde

\[
 \mathcal H_{P,\mathfrak M}^{\rm surv}
 =\Pi_{P,\mathfrak M}^{\rm surv}\mathcal H_F,
\]
\[
 \widehat{\mathcal M}_{P,\mathfrak M}
 =\Pi_{P,\mathfrak M}^{\rm surv}
  \widehat{\mathcal M}_F
  \Pi_{P,\mathfrak M}^{\rm surv},
\]
y \(\mu_{P,\mathfrak M}\) es su medida espectral en el estado preparado.

Cuando el acoplamiento se especifica por una representación irreducible \(\lambda\) de un grupo finito \(G_P\), el proyector correspondiente se obtiene sin elegir una ruta distinguida:

\[
 P_{P,\lambda}
 =\frac{\dim\lambda}{|G_P|}
  \sum_{g\in G_P}\overline{\chi_\lambda(g)}\,U(g).
\]
La salida bidireccional es entonces el par comprimido

\[
 \mathbb M(P,\lambda)=
 \left(
 P_{P,\lambda}\widehat{\mathcal M}^{\beta,D}P_{P,\lambda},
 P_{P,\lambda}\widehat{\mathcal M}^{\beta,\Omega}P_{P,\lambda}
 \right).
\]
Si ambos operadores conmutan con la representación irreducible, el lema de Schur fuerza su escalaridad sobre el bloque. Si no conmutan, el resultado correcto es el espectro del bloque mínimo invariante.

\Needspace{7\baselineskip}
\subsection{Naturalidad del funtor espectral bajo entrelazadores de la dinámica TPK, el acoplamiento y el operador de masa}
Si \(V:F\to F'\) entrelaza la dinámica TPK, el núcleo de acoplamiento y los operadores de masa, entonces

\[
 V\Pi_{P,\mathfrak M}^{\rm surv}
 =\Pi_{P',\mathfrak M'}^{\rm surv}V,
 \qquad
 V\widehat{\mathcal M}_{P,\mathfrak M}
 =\widehat{\mathcal M}_{P',\mathfrak M'}V.
\]
Por el cálculo funcional, la medida espectral se transporta por \(V\). La salida no depende de una numeración de las celdas ni de la elección de una base dentro de una órbita. Esta es la condición que debe sustituir a todo selector nominal no equivariante.

\Needspace{7\baselineskip}
\subsection{Condiciones de existencia de una masa escalar bajo el lector dimensional}
Una masa escalar está determinada en cualquiera de los casos siguientes:

\begin{enumerate}
\item el subespacio persistente tiene rango uno;
\item una representación irreducible fija un único proyector espectral;
\item el protocolo de medida define un estado positivo normalizado
\(\omega_{P,\mathfrak M}\) sobre el álgebra de observables;

\item una resonancia define un polo aislado del resolvente.
\end{enumerate}
En el tercer caso,

\[
 m(P,\mathfrak M)
 =\omega_{P,\mathfrak M}
  (\widehat{\mathcal M}_{P,\mathfrak M}).
\]
En ausencia de uno de estos mecanismos, la salida científica es el operador, su espectro y sus multiplicidades. Esa salida no es incompleta: es el objeto que la dinámica ha determinado.

\Needspace{7\baselineskip}
\subsection{Construcción finita del núcleo de masa mediante filtros de representación, persistencia y acoplamiento, seguida de normalización irreducible}
El núcleo se obtiene en cuatro etapas finitas:

\begin{enumerate}
\item \textbf{Filtro de representación.} Se retienen las rutas cuyo sector, rama de carga,
color, holonomía \(2\pi/4\pi\), paridad y estadística coinciden con \(P\).

\item \textbf{Filtro de persistencia.} Se usa el orden parcial formado por retorno de clase,
retorno de órbita, retorno de celda, separación de prefijos, deuda de acarreo y estabilidad decimal. El resultado puede ser una órbita de Pareto; no se fuerza unicidad.

\item \textbf{Filtro de acoplamiento.} Se evalúa la incidencia de frontera y el canal de
interacción. Rutas desacopladas reciben peso cero.

\item \textbf{Normalización.} Sobre cada componente irreducible se normaliza la traza del
núcleo a uno.

\end{enumerate}
El atlas orientado proporciona todos los datos de las dos primeras etapas. La tercera exige que cada realización física declare su canal; la cuarta es canónica una vez fijadas las componentes irreducibles.

\Needspace{7\baselineskip}
\section{Transductor genealógico de torres}
\Needspace{7\baselineskip}
\subsection{Par de palabras asociadas a una misma ruta genealógica}
Cada \(\Gamma\in\mathcal R_{324}\) determina simultáneamente una palabra temporal

\[
 w_\Gamma\in\mathbb F_3^{108}
\]
y una palabra dodecafásica firmada

\[
 \tau_\Gamma\in\{-1,0,+1\}^{12}.
\]
La primera conserva desplazamientos a lo largo de CT--108; la segunda conserva la orientación acumulada de los doce acontecimientos. Las dos palabras están ligadas por el mismo registro genealógico, pero no son intercambiables. El refinamiento de torres lee ambas y conserva, además, hoja, acarreo, devanado y memoria de retorno.

\Needspace{7\baselineskip}
\subsection{Correspondencia entre altura armónica y longitud temporal}
Escribamos

\[
 \mathfrak S=30\langle3,4\rangle
 =\{90,120\}\cup\{30r:r\ge6\}.
\]
Para \(h=30r\), la longitud temporal asociada es

\[
 \ell_h=9r.
\]
En particular,

\[
 90\longleftrightarrow27,\qquad
 120\longleftrightarrow36,\qquad
 360\longleftrightarrow108.
\]
Esta correspondencia no identifica la torre con el tiempo: entrelaza la altura del operador armónico con la longitud del segmento de ruta que la genera.

\Needspace{7\baselineskip}
\subsection{Lector de desplazamiento del ciclo temporal de 108 estados}
Definimos la discrepancia cíclica

\[
 D_\Gamma(n)=
 \sum_{t=0}^{107}
 \mathbf1_{\{w_{\Gamma,t}\ne
 w_{\Gamma,t+n\!\!\!\pmod{108}}\}}.
\]
El coeficiente relativo de la rama temporal es

\[
 \boxed{\eta_{30r}^{D}(\Gamma)
 =D_\Gamma(9r)-D_{\Gamma_e}(9r).}
\]
La resta no usa una masa como origen: usa la ruta central como origen afín del carácter relativo. Se cumple

\[
 \eta_{90}^{D}+\eta_{120}^{D}
 =[D_\Gamma(27)+D_\Gamma(36)]-80,
\]
de modo que la primera capa prolonga exactamente el primer componente del lector \(K_D\).

\Needspace{7\baselineskip}
\subsection{Lector de acumulación dodecafásica}
Para \(r\ge1\), sea

\[
 A_\Gamma(r)=
 \sum_{j=0}^{11}
 \left(
 \sum_{u=0}^{r-1}
 \tau_{\Gamma,j+u\!\!\!\pmod{12}}
 \right)^2.
\]
El coeficiente relativo de la rama orientada es

\[
 \boxed{\eta_{30r}^{\Omega}(\Gamma)
 =A_\Gamma(r)-A_{\Gamma_e}(r).}
\]
Entonces

\[
 4\eta_{90}^{\Omega}-3\eta_{120}^{\Omega}
 =\Omega(\tau_\Gamma)-80,
\]
por lo que esta capa prolonga el primer componente de \(K_\Omega\). Las dos familias forman el levantamiento bidireccional

\[
 \mathcal F_{\rm tower}^{2w}:
 \mathcal R_{324}\longrightarrow
 \mathcal E_{\mathfrak S}\times\mathcal E_{\mathfrak S},
 \qquad
 \Gamma\longmapsto
 \bigl(\eta^D(\Gamma),\eta^\Omega(\Gamma)\bigr).
\]
\Needspace{7\baselineskip}
\subsection{Determinación finita de la torre infinita}
La periodicidad de la palabra temporal da

\[
 D_\Gamma(9(r+12))=D_\Gamma(9r).
\]
Si

\[
 T_\Gamma=\sum_{j=0}^{11}\tau_{\Gamma,j},
 \qquad r=12q+s,\quad 0\le s<12,
\]
entonces

\[
 A_\Gamma(12q+s)
 =A_\Gamma(s)+(12q^2+2qs)T_\Gamma^2.
\]
Por tanto, toda la torre queda determinada por doce valores de \(D_\Gamma\), doce valores de \(A_\Gamma\) y el entero \(T_\Gamma^2\). La infinitud de alturas no exige una infinitud de datos independientes.

La misma fórmula resuelve el diamante genealógico de altura 360. Si

\[
 h(a,b)=90a+120b,
\]
entonces

\[
 (a+4,b)\sim(a,b+3)
\]
porque \(3(a+4)+4b=3a+4(b+3)\). La genealogía \((a,b)\) se conserva hasta aplicar memoria y acarreo; sólo después desciende a la altura común.

\Needspace{7\baselineskip}
\subsection{Convergencia absoluta del carácter refinado sobre las torres \(90/120\)}
Se tiene

\[
 |\eta_h^D|\le108,\qquad
 |\eta_{30r}^{\Omega}|\le24r^2.
\]
Como \(0<q_+<q_-<1\),

\[
 |s_{30r}|\le2q_-^{30r}.
\]
De aquí se deduce

\[
 \sum_{h\in\mathfrak S}|\eta_h^D s_h|<\infty,
 \qquad
 \sum_{h\in\mathfrak S}|\eta_h^\Omega s_h|<\infty.
\]
El carácter refinado está bien definido sin truncar la torre.

\Needspace{7\baselineskip}
\subsection{Memoria, acarreo y cociclo}
El coeficiente desnudo se enriquece mediante

\[
 \mathbb F_{\rm tower}(\Gamma)_h
 =\bigl(\eta_h^D,\eta_h^\Omega,\rho_h,
 \mathbf w_h,c_h,N_h^{\rm surv}\bigr),
\]
donde la firma truncada satisface

\[
 \sigma_{\Gamma|\ell_h}
 =r(\rho_h)+\mathcal D\mathbf w_h,
 \qquad
 \mathcal D=\operatorname{diag}(2,6,2,2,4),
\]
y el acarreo de dos residuos es

\[
 c(\rho_1,\rho_2)=
 \mathcal D^{-1}
 \bigl[r(\rho_1)+r(\rho_2)-r(\rho_1+\rho_2)\bigr].
\]
El defecto de localidad de una composición,

\[
 \delta_h(\Gamma,\Lambda)=
 F_h(\Gamma\star\Lambda)-F_h(\Gamma)-F_h(\Lambda),
\]
obedece la identidad de cociclo

\[
 \delta_h(\Gamma,\Lambda)
 +\delta_h(\Gamma\star\Lambda,\Xi)
 =\delta_h(\Lambda,\Xi)
 +\delta_h(\Gamma,\Lambda\star\Xi).
\]
Así se preserva por qué dos rutas con la misma altura visible pueden transportar historias distintas.

\Needspace{7\baselineskip}
\subsection{Selección de la lectura por el acoplamiento}
Las dos ramas no se promedian por conveniencia. El acoplamiento físico \(a\) debe determinar una transformación natural

\[
 \Pi_a:
 \mathcal E_{\mathfrak S}\times\mathcal E_{\mathfrak S}
 \longrightarrow\mathcal E_{\mathfrak S}
\]
que conmute con truncaciones y simetrías, respete el cociclo de acarreo, conserve el origen electrónico y sea compatible con \(X_{90}^{4}=Y_{120}^{3}\). Si el acoplamiento intercambia las dos lecturas, la única clausura continua, multiplicativa, simétrica y normalizada sobre escalares positivos es la media geométrica; en exponentes,

\[
 L_{\rm sym}=\frac{L_D+L_\Omega}{2}.
\]
En ausencia de esa simetría, la salida es la medida espectral conjunta de \((L_D,L_\Omega)\), no una media arbitraria.

\Needspace{7\baselineskip}
\subsection{Refinamiento coinductivo por órbitas primitivas}
El levantamiento finito anterior admite una prolongación. Sea \(\mathscr X_\Gamma\) el grafo de extensiones admisibles de la ruta, \(U_\Gamma\) su operador de transferencia y \(R_\Gamma\) la involución de orientación. Las trazas

\[
 a_n(\Gamma)=\operatorname{Tr}(R_\Gamma U_\Gamma^n)
\]
y la inversión de Möbius

\[
 p_n(\Gamma)=\frac1n\sum_{d\mid n}\mu(d)a_{n/d}(\Gamma)
\]
separan órbitas primitivas. La diferencia

\[
 p_h(\Gamma)-p_h(\Gamma_e)
\]
es un refinamiento ulterior del coeficiente de altura \(h\), no una sustitución de las dos lecturas finitas.

\Needspace{7\baselineskip}
\subsection{Formulación histórica mediante autómata de extensión}
La firma \(\mathbb Z^5\) no agota la historia de una ruta. Para cada \(\Gamma\in\mathcal R_{324}\) se conservan:

\begin{itemize}
\item las 108 transiciones del registro genealógico;
\item los veinte bloques y sus prefijos;
\item las palabras \(w_6\);
\item las nueve comparaciones \(K\leftrightarrow K+9\);
\item la hoja, el acarreo y el devanado;
\item los primeros retornos de clase, órbita, celda y estado cinemático;
\item la ambigüedad acumulada y las tríadas decimales estabilizadas;
\item los perfiles \(K_D\) y \(K_\Omega\).
\end{itemize}
Estos datos permiten construir un transductor de torres sin abrir el residual de una masa.

\Needspace{7\baselineskip}
\subsection{Autómata local de extensión}
Sea \(\mathscr X_\Gamma\) el conjunto finito de estados locales compatibles con la ruta. Un estado conserva

\[
 x=(t,g,B,p,\varepsilon,c,w),
\]
donde \(t\) es el tiempo discreto, \(g\) la fase nonádica, \(B\) el bloque, \(p\) el prefijo, \(\varepsilon\) la hoja, \(c\) el acarreo y \(w\) el devanado. Existe una arista \(x\to y\) cuando \(y\) es una extensión admisible de \(x\), respeta el registro genealógico y supera la poda \(G_9\). La tabla de 36.864 composiciones de acarreo proporciona la ley local de suma; las tablas de retorno fijan la memoria.

Sea \(U_\Gamma\) el operador de transferencia normalizado de este grafo y \(R_\Gamma\) la involución de orientación. Se exige

\[
 U_{C_M\Gamma}=J U_\Gamma J^{-1},\qquad
 R_{C_M\Gamma}=-J R_\Gamma J^{-1}.
\]
\Needspace{7\baselineskip}
\subsection{Conteos orientados de celdas primitivas y construcción de la firma electrónica}
Para \(n\ge1\), la traza orientada

\[
 a_n(\Gamma)=\operatorname{Tr}(R_\Gamma U_\Gamma^n)
\]
cuenta retornos cerrados con signo. La inversión de Möbius separa órbitas primitivas:

\[
 p_n(\Gamma)
 =\frac1n\sum_{d\mid n}\mu(d)a_{n/d}(\Gamma).
\]
Se toma el electrón como origen afín de las correcciones relativas y se define

\[
 \boxed{
 \eta_h(\Gamma)=p_h(\Gamma)-p_h(\Gamma_e),
 \qquad h\in\mathfrak S.}
\]
Así se obtiene el mapa

\[
 \mathcal F_{\rm tower}:
 \Gamma^{\rm enr}\longmapsto
 \eta(\Gamma)=(\eta_h(\Gamma))_{h\in\mathfrak S}.
\]
No hay parámetros por especie. El mismo autómata, la misma involución y la misma inversión de Möbius actúan sobre las 324 rutas.

\Needspace{7\baselineskip}
\subsection{Convergencia del refinamiento primitivo}
Sea \(N_\Gamma=\dim\mathscr X_\Gamma\). Si \(U_\Gamma\) es una contracción,

\[
 |a_n(\Gamma)|\le N_\Gamma.
\]
El número de divisores de \(n\) crece sublinealmente, por lo que

\[
 |p_n(\Gamma)|
 \le\frac{N_\Gamma\tau(n)}n.
\]
Como \(0<q_+<q_-<1\),

\[
 |s_h|\le q_-^h+q_+^h\le2q_-^h.
\]
Por comparación,

\[
 \sum_{h\in\mathfrak S}|\eta_h(\Gamma)s_h|<\infty
\]
uniformemente en cualquier colección finita de fibras. El carácter refinado está, por tanto, bien definido.

\Needspace{7\baselineskip}
\subsection{Identidades de conservación de la ruta electrónica y unicidad del punto fijo central}
La construcción debe satisfacer simultáneamente:

\begin{enumerate}
\item invariancia por renumeración de estados;
\item covarianza bajo conjugación de hoja;
\item compatibilidad con suma directa de componentes;
\item igualdad al recomputar desde registro genealógico y desde el grafo certificado;
\item conservación del origen \(\eta(\Gamma_e)=0\);
\item insensibilidad absoluta a una permutación de las columnas externas de masa;
\item determinación de las colas mediante una cota anterior al contraste.
\end{enumerate}
La representación exacta de un residual dado demuestra completitud del codominio, pero no sustituye estas siete pruebas.

\Needspace{7\baselineskip}
\section{Reloj de ruta y normalización absoluta}
\Needspace{7\baselineskip}
\subsection{Los retornos enteros no bastan para distinguir especies}
En el atlas de 324 rutas, el primer retorno de celda es 27, el retorno cinemático dentro de CT--108 es 54 y el retorno cinemático extendido es 216. Al ser comunes, estos enteros fijan la arquitectura temporal global, pero no pueden explicar por sí solos diferencias de masa entre especies.

\Needspace{7\baselineskip}
\subsection{Holonomía temporal de la ruta electrónica y retorno con memoria}
La información específica reside en la holonomía del operador de extensión. Sea \(\lambda_j(\Gamma)\) un autovalor no trivial de \(U_\Gamma\) sobre el subespacio persistente. Escribimos

\[
 \lambda_j(\Gamma)=r_j(\Gamma)e^{i\theta_j(\Gamma)}.
\]
La frecuencia interna asociada es

\[
 \omega_j(\Gamma)
 =\frac{\theta_j(\Gamma)+2\pi w_j(\Gamma)}{108t_0},
\]
donde el entero \(w_j\) lo fija el devanado registrado, no una elección de rama posterior. Para un modo estable \(r_j=1\); para una resonancia \(r_j<1\) y su atenuación aporta una anchura.

El cociente

\[
 \rho_{\omega,j}(\Gamma)
 =\frac{\omega_j(\Gamma)}{\omega_0}
\]
es adimensional y puede modificar razones de masa sin alterar la década común \(-34\). La energía y la masa del modo son

\[
 E_j(\Gamma)=\hbar_{\rm ret}\omega_j(\Gamma)
 \chi_{\rm full}(\Gamma)R_{\rm act}^{\kappa_j(\Gamma)},
\]
\[
 m_j(\Gamma)=E_j(\Gamma)c_{\rm int}^{-2}.
\]
Esta ecuación reúne acción, reloj, carácter y lector sin convertir ninguno de ellos en otro. Su aplicación numérica exige determinar \(U_\Gamma\), fijar el devanado y declarar la carta dimensional.

\Needspace{7\baselineskip}

